% Generated draft: requires mathematical and source review.
\documentclass[11pt]{article}
\usepackage[margin=25mm]{geometry}
\usepackage{fontspec}
\setmainfont{Latin Modern Roman}
\usepackage{amsmath,amssymb,mathtools,unicode-math}
\setmathfont{Latin Modern Math}
\usepackage{enumitem,needspace,xurl,hyperref}
\setlength{\parindent}{0pt}
\setlength{\parskip}{5pt}
\setlength{\emergencystretch}{3em}
\newcommand{\N}{\mathbb{N}}
\newcommand{\Z}{\mathbb{Z}}
\newcommand{\Q}{\mathbb{Q}}
\newcommand{\R}{\mathbb{R}}
\newcommand{\C}{\mathbb{C}}
\newcommand{\F}{\mathbb{F}}
\newcommand{\A}{\mathbb{A}}
\newcommand{\PP}{\mathbb{P}}
\newcommand{\E}{\mathbb{E}}
\newcommand{\eps}{\varepsilon}
\newcommand{\dd}{\,\mathrm{d}}
\begin{document}
\section*{UnsolvedMath ID: 30002711}
\textbf{Source problem code:} catalogue entry
\subsection{Definitions and mathematical statement}
The supplied question asks whether the local lifting problem for cyclic groups \(G\) outside an unspecified known construction can always be solved with
\[
R=W(k)(\zeta_{|G|}),
\]
where \(\zeta_{|G|}\) denotes a primitive root of unity of order \(|G|\). The record does not specify the hypotheses on \(k\), the groups covered by the known construction, or the precise meaning of this coefficient-ring notation. Consequently, it does not determine a fully specified mathematical assertion.

The following describes the usual local lifting formulation, conditionally on the missing conventions. Suppose that \(k\) is a perfect field of characteristic \(p>0\), so that \(W(k)\) denotes its ring of \(p\)-typical Witt vectors. Write \(k[[t]]\) for the formal power series ring. A faithful local action of a finite cyclic group \(G\) is an injective homomorphism
\[
\rho:G\longrightarrow\operatorname{Aut}^{\mathrm{cont}}_k(k[[t]]),
\]
where continuity is for the \((t)\)-adic topology. For a complete discrete valuation ring \(R\) of characteristic zero with maximal ideal \(\mathfrak m_R\) and a specified identification \(R/\mathfrak m_R\cong k\), a lift of \(\rho\) is an injective homomorphism
\[
\widetilde\rho:G\longrightarrow\operatorname{Aut}^{\mathrm{cont}}_R(R[[T]])
\]
whose reduction modulo \(\mathfrak m_R\), under \(T\mapsto t\), equals \(\rho\). Continuity here is for the \((\mathfrak m_R,T)\)-adic topology.

Under this conventional interpretation, the intended question would ask whether, for every cyclic group outside the groups covered by the referenced construction and every faithful local action under consideration, such a lift exists over the prescribed ring \(W(k)(\zeta_{|G|})\). This universal quantification is an interpretation of ``the local lifting problem''; the supplied record does not explicitly state its scope or the allowed fields \(k\).

The prescribed ring requires clarification: adjoining \(\zeta_{|G|}\) integrally to \(W(k)\), adjoining it to the fraction field, and taking a suitable ring of integers are distinct descriptions. A coefficient ring for the formulation above must also have residue field identified with \(k\). No choice among these interpretations, no additional hypothesis on \(k\), and no description of the excluded known construction is supplied here. These qualifications remain unresolved.
\subsection{Short English statement}
The question asks whether cyclic local actions outside an unspecified known construction lift to characteristic zero over the particular coefficient ring written \(W(k)(\zeta_{|G|})\). The field hypotheses, excluded cases, and exact coefficient ring need clarification. The supplied dated assessment labels the stronger coefficient-ring question open in wild cyclic cases; this draft does not independently verify that assessment.
\subsection{Sources}
\begin{itemize}
\item Local UnsolvedMath registry: supplied complete record for permanent ID \(30002711\), title ``Minimal Coefficient Rings for Cyclic Local Lifts,'' exact source problem code \texttt{catalogue entry}, in the Oberwolfach Reports Open Problems collection. The supplied proposed year is 2014; no individual proposer is identified.
\item Supplied report citation: \emph{Valuation Theory and Its Applications} (2015), Oberwolfach Reports, DOI 10.4171/owr/2014/49, pp.\ 2757--2823. \url{https://doi.org/10.4171/owr/2014/49}. The record describes this as a workshop report containing open-problem discussions.
\item Supplied literature source: ``Oberwolfach Report 49/2014.'' \url{https://publications.mfo.de/bitstream/handle/mfo/3439/OWR_2014_49.pdf?isAllowed=y&sequence=1}.
\item Supplied literature source: Obus, \emph{A generalization of the Oort conjecture}. \url{https://arxiv.org/abs/1502.07623}.
\item Supplied literature assessment, checked on exactly \texttt{2026-08-21}: existence of a cyclic lift is reported as known, while the stronger assertion prescribing \(W(k)\) with the indicated roots of unity is reported as open for wild cyclic cases. This is supplied dated curation, described in the record as literature triage derived from the supplied OWR curation; it has not been independently verified here. The record requests independent verification of time-sensitive status and cited claims.
\item The registry labels the problem \texttt{open} and the statement \texttt{corrected\_verified}, and reports that it was checked against the official report. These are supplied registry assertions, not verification performed for this draft.
\end{itemize}
\end{document}
